\chapter{GAME SYSTEM}

% ============================================================================
% Section 5.1: Software Overview
% ============================================================================
\section{Software Overview}

\textit{Wiz's Journey: Mastering the States of Matter} is developed using Unity 2022 LTS, a cross-platform game engine selected for its extensive capabilities in 2D game development, robust documentation, and the development team's prior expertise. The game is designed as a 2D platformer that integrates educational content with interactive gameplay mechanics, allowing students to progress through levels while encountering learning outcomes related to phase changes of matter.

The game is deployed across multiple platforms to ensure widespread accessibility and reach diverse student populations. The primary deployment target is Android mobile devices, distributed through Itch.io, a popular game distribution platform that allows easy access without requiring app store approval processes. A secondary web-based version is deployed via Unity Play, allowing students to play directly in web browsers on computers and tablets without installation requirements. This multi-platform approach reflects the varying technological infrastructure available to schools and students in the Philippines and aligns with the goal of making the educational game accessible regardless of device availability.

The development timeline began in December 2025 at the conclusion of the first trimester and extended through March 2026, spanning the academic term break and into the second trimester. This timeline was established to accommodate the academic calendar constraints and allowed the development team to complete the game while managing other coursework obligations. By the time development concluded in March 2026, the target Grade 6 students had already completed their first quarter instruction on phase changes under the MATATAG curriculum, positioning the game as a supplementary learning reinforcement tool rather than a primary instructional resource.

The game architecture incorporates both lower-order and higher-order mechanics aligned with the selected learning outcomes. Lower-order mechanics focus on foundational understanding of the four phase changes—melting, freezing, evaporation, and condensation—using visual representations and simple interactions. Higher-order mechanics build upon this foundation to address more complex concepts such as the relationship between heat transfer and state changes, the particulate nature of matter during phase transitions, and critical thinking about real-world applications of phase changes. The system integrates a particle lattice visualization system to support students' conceptualization of microscopic-level processes underlying macroscopic observations, addressing a common misconception identified in science education research.
